\documentclass{aastex701}

\begin{document}

\title{Five Transiting Planet Candidates from the TESS Quick-Look Pipeline Light Curves of Sectors 101 and 103}

\author[0000-0000-0000-0000]{Carl Drinkwater}
\affiliation{Independent researcher, Shropshire, UK}
\email{carl@29degrees.co.uk}

\begin{abstract}
I report five transit signals found in a search of the TESS Quick-Look Pipeline light curves of Sectors 101 and
103 for stars brighter than $T=11.5$. Each signal repeats in two to six other TESS sectors, falls on the target in
difference images, shows no secondary eclipse and no significant odd-even depth difference, and matches no TOI,
CTOI, SPOC threshold crossing event or TESS eclipsing binary. Joint fits of all sectors give radii of 2.6 to
11 $R_\oplus$ and periods of 0.76 to 3.96 d, with TRICERATOPS false positive probabilities of 0.010 to 0.47. Two are
ultra-short-period candidates. Three hosts show Gaia radial velocity variability or a high RUWE, which needs
follow-up.
\end{abstract}

\section{Search and vetting}

I searched the Quick-Look Pipeline \citep[QLP;][]{huang2020a,huang2020b,kunimoto2021} light curves of TESS
\citep{ricker2015} Sectors 101 and 103 for the stars brighter than $T=11.5$: 298,298 and 277,524 stars. I removed
trends with a biweight filter of 0.5 d \citep{hippke2019} and searched 0.5 to 14 d with box least squares
\citep{kovacs2002}. I kept the signals that pass every flux-level false alarm and false positive test of LEO-vetter
\citep{kunimoto2025} and match no TOI \citep{guerrero2021}, CTOI, SPOC threshold crossing event or TESS eclipsing
binary \citep{prsa2022} in the lists of 2026 November 7.

For each signal I searched the TESS-SPOC \citep{caldwell2020} and QLP light curves of the other sectors of the
star. A joint box least squares search of all sectors at the candidate period had to reach 1.5 times the highest
SNR at ten null periods of 0.71 to 1.41 times that period. I fitted a point spread function model to the in-transit difference images \citep{bryson2013} of the TESScut
cutouts \citep{brasseur2019} of each sector, and of all sectors together, to locate the source of the dip among the
Gaia DR3 \citep{gaia2023} stars. I then fitted one transit model \citep{kreidberg2015} to all sectors with emcee
\citep{foremanmackey2013}, measured the odd and even depths and the flux at phase 0.5, and ran TRICERATOPS
\citep{giacalone2021} five times on the folded light curve, with the nearby-star scenarios removed where the
difference images place the dip on the target. The radii use the TIC stellar radius \citep{stassun2019}.

Table~\ref{tab:candidates} lists the five candidates. For each, the joint search SNR is 18 to 47 against null maxima
of 5.5 to 15.4, the odd and even depths differ by less than 3$\sigma$, and the flux at phase 0.5 is within
1$\sigma$ of zero.

\begin{table*}[ht!]
\centering
\caption{Transit parameters from the joint fit of all TESS sectors. $T_0$ is in BJD$_\mathrm{TDB} - 2457000$.
Uncertainties in the last digits are in parentheses. FPP is the mean TRICERATOPS false positive probability.
\label{tab:candidates}}
\begin{tabular}{rcllrrrrc}
\hline
TIC & $T$ & $P$ (d) & $T_0$ & Depth (ppm) & Dur.\ (h) & $R_p$ ($R_\oplus$) & FPP & Sectors \\
\hline
334982989 & 11.23 & 0.7646381(6) & 4069.6852(7) & $707\pm36$ & $1.30\pm0.05$ & $2.65\pm0.19$ & 0.017 & 10, 64, 99--101 \\
273213461 & 11.15 & 0.8565034(14) & 4069.8242(14) & $554\pm29$ & $2.81\pm0.09$ & $4.54\pm0.35$ & 0.010 & 10, 64, 99--101 \\
242415665 & 11.05 & 2.9344279(23) & 4104.9671(11) & $1948\pm58$ & $1.95\pm0.07$ & $5.61\pm0.49$ & 0.47 & 11, 38, 65, 101, 102 \\
437823249 & 9.87 & 3.2223295(66) & 3093.7904(16) & $627\pm26$ & $3.60\pm0.13$ & $8.0\pm0.8$ & 0.030 & 11, 38, 64, 65, 99--101 \\
28812587 & 11.07 & 3.9634883(25) & 4156.1382(12) & $5807\pm187$ & $2.06\pm0.10$ & $11.2\pm1.1$ & 0.11 & 12, 39, 103 \\
\hline
\end{tabular}
\end{table*}

\section{Notes on individual candidates}

\textbf{TIC 334982989} is an ultra-short-period sub-Neptune candidate around a G dwarf of 0.97 $R_\odot$ at 214 pc.
The difference images of Sectors 100 and 101 place the dip on the target (joint SNR 11.5). Gaia DR3 gives a RUWE
of 0.99 and no radial velocity variability in 27 epochs \citep{katz2023}, which limits a companion on the transit
orbit to below 11 Jupiter masses.

\textbf{TIC 273213461} is an ultra-short-period Neptune-size candidate around an F star of 1.94 $R_\odot$.
TRICERATOPS gives the lowest FPP of the five, but Gaia DR3 gives a RUWE of 3.01 and radial velocity variability with
an epoch scatter of 10.3 km\,s$^{-1}$. A close stellar companion is likely; its light would dilute the transit, so
the true radius can be larger.

\textbf{TIC 242415665} recurs in five sectors with a joint search SNR of 29.7. Its odd and even depths differ by
2.98$\sigma$ and its FPP is 0.47, so an eclipsing binary at twice the period remains possible. The fitted stellar
density is 6.7 times the TIC value, which suggests a smaller star and companion than the table gives.

\textbf{TIC 437823249} has a fitted stellar density 14 times the TIC value. The TIC gives a subgiant of
3.55 $R_\odot$, but the transit fits a dwarf of about 1.5 $R_\odot$, for which the companion would have a radius of
3 to 4 $R_\oplus$. Gaia DR3 shows radial velocity variability with an epoch scatter of 3.0 km\,s$^{-1}$.

\textbf{TIC 28812587} is a giant planet candidate around a K star at 200 pc. The Sector 103 difference image places
the dip on the target (SNR 16.1). Gaia DR3 shows radial velocity variability (14 epochs, scatter
5.5 km\,s$^{-1}$), which limits a companion on the transit orbit to below 0.051 $M_\odot$; a third body or Gaia
noise can also cause it. The fitted stellar density is 5.0 times the TIC value, and the TIC contamination ratio is
0.58.

\section{Conclusion}

These five signals are planet candidates, not validated planets. Radial velocities can measure the masses and test
the three hosts with Gaia radial velocity variability or a high RUWE for stellar companions, and high-resolution
imaging can test for unresolved neighbours. I will upload the five to ExoFOP as Community TOIs.

\begin{acknowledgments}
This note uses data from the TESS mission, obtained from MAST. Claude (Anthropic) wrote the analysis code and a draft of
this note under my direction. The code and results are at https://github.com/carldr/tess-transiting-planet-search.
\end{acknowledgments}

\facilities{TESS, Gaia}

\software{astropy, astroquery, batman \citep{kreidberg2015}, emcee \citep{foremanmackey2013}, LEO-vetter
\citep{kunimoto2025}, lightkurve, TRICERATOPS \citep{giacalone2021}, wotan \citep{hippke2019}}

\bibliography{references}{}
\bibliographystyle{aasjournalv7}

\end{document}
